\chapter{Exercício 6}
\label{cha:exercise_6}

\section{Assinatura digital com RSA}

O exercício 6 propõe a implementação de um esquema de assinatura digital baseado no algoritmo assimétrico \textbf{RSA}, usando a biblioteca \textit{Java Cryptography Architecture (JCA)}.  
A aplicação permite assinar e verificar ficheiros de qualquer dimensão, suportando as funções de hash \texttt{SHA-1} e \texttt{SHA-256}, e valida a cadeia de certificação do certificado do assinante antes de verificar a assinatura.

\subsection{Descrição geral da implementação}

A aplicação \texttt{DigitalSign.java} define duas operações principais configuradas por linha de comandos: o modo de assinatura (\texttt{-sign}) e o modo de verificação (\texttt{-verify}).  
O algoritmo de assinatura selecionado depende da flag de hash: \texttt{SHA1withRSA} para \texttt{-sha1} e \texttt{SHA256withRSA} para \texttt{-sha256}.

A implementação faz uso das APIs padrão de segurança do Java para:
\begin{itemize}
    \item Carregar \textit{keystores} (formato PKCS\#12) e extrair chaves privadas;
    \item Criar e utilizar objetos \texttt{Signature} para assinar/verificar dados;
    \item Carregar certificados X.509 e construir/validar caminhos de certificação com \texttt{CertPathBuilder} e parâmetros PKIX.
\end{itemize}

\subsection{Modo de assinatura (\texttt{-sign})}

No modo de assinatura, o programa executa os seguintes passos (descritos conceptualmente, sem apresentar código):

\begin{enumerate}
    \item O \textit{keystore} em formato PKCS\#12 é carregado utilizando a palavra-passe fornecida pelo utilizador.
    \item O programa seleciona automaticamente a primeira entrada de chave privada encontrada no keystore. É feita uma verificação para garantir que a entrada obtida é efectivamente uma chave privada.
    \item É inicializado um objecto de assinatura com o algoritmo apropriado (SHA-1 ou SHA-256 com RSA).
    \item O ficheiro a assinar é lido de forma incremental, em blocos de dimensão fixa (no código de referência usa-se um buffer de 8\,192 bytes), e cada bloco é fornecido ao objecto de assinatura para atualizar o seu estado de forma incremental.
    \item Quando todo o ficheiro foi processado, é gerada a assinatura final (um array de bytes) que é gravada num ficheiro de saída dedicado.
\end{enumerate}

A leitura incremental e a actualização por blocos do objecto de assinatura permitem que ficheiros de tamanho arbitrário sejam assinados sem carregar o conteúdo inteiro na memória.

\subsection{Modo de verificação (\texttt{-verify})}

No modo de verificação, a aplicação procede conceptualmente assim:

\begin{enumerate}
    \item Carrega o certificado do assinante (formato X.509).
    \item Carrega o \textit{truststore} (formato JKS) que contém as entidades certificadoras consideradas confiáveis.
    \item Carrega certificados intermediários a partir de uma directoria configurada, para auxiliar na construção da cadeia.
    \item Constrói e valida a cadeia de certificação do certificado do assinante usando o \texttt{CertPathBuilder} e parâmetros PKIX. Caso a cadeia não seja válida, o processo de verificação da assinatura é interrompido e sinaliza-se a falha de confiança.
    \item Se a cadeia for válida, inicializa-se o objecto de verificação com a chave pública do certificado do assinante.
    \item O ficheiro cujo conteúdo se pretende verificar é lido incrementalmente (blocos de 8\,192 bytes) e cada bloco é fornecido ao objecto de verificação para actualizar o estado. No fim, a assinatura armazenada é comparada com o resultado calculado para determinar se a assinatura é válida.
\end{enumerate}

A validação da cadeia antes da verificação assegura que apenas certificados confiáveis são usados para aceitar assinaturas.

\subsection{Processamento de ficheiros de grande dimensão}

A característica fundamental que permite ao programa processar ficheiros de qualquer dimensão é o processamento em streaming (leitura incremental com buffer):

\begin{itemize}
    \item Em vez de carregar todo o ficheiro na memória, o programa lê blocos de tamanho fixo (no código de referência: 8\,192 bytes) e processa cada bloco imediatamente.
    \item Nas operações de assinatura e de verificação, os métodos \texttt{update()} do objecto \texttt{Signature} são chamados por cada bloco lido, permitindo que o cálculo do digest/assinatura seja efectuado de forma incremental.
    \item Só é mantido em memória um pequeno buffer por cada iteração; o resto do ficheiro permanece em disco até ser lido. Isto garante baixa utilização de memória mesmo para ficheiros muito grandes.
\end{itemize}

Esta abordagem torna a aplicação escalável e apropriada para ficheiros desde alguns kilobytes até gigabytes, dependendo apenas das capacidades de I/O do sistema e não da RAM disponível.

\subsection{Validação da cadeia de certificação}

A validação da cadeia é tratada através das seguintes decisões de design (descritas conceptualmente):

\begin{itemize}
    \item O certificado do assinante é utilizado como o ponto de partida (end-entity certificate).
    \item Certificados intermédios são adicionados a um repositório temporário que auxilia a construção do caminho até a(s) raiz(es) de confiança presentes no \textit{truststore}.
    \item O mecanismo de construção e validação usa o modelo PKIX: são definidos parâmetros que incluem o \textit{truststore} como fonte de raízes confiáveis e o certificado do assinante como alvo, e é invocado um \texttt{CertPathBuilder} para tentar construir e validar um caminho válido.
    \item Caso a construção falhe (cadeia incompleta, certificados expirados, ou outros erros de validação), o programa rejeita o certificado e reporta a cadeia como inválida, evitando assim verificar assinaturas com certificados não confiáveis.
\end{itemize}

\subsection{Execução e utilização}

A aplicação é controlada por argumentos na linha de comandos, distinguindo o modo de assinatura do modo de verificação e a função de hash a empregar. Em termos gerais, o fluxo de operação é:

\begin{itemize}
    \item No modo de assinatura: fornecer ficheiro a assinar, \textit{keystore}, palavra-passe e o ficheiro onde guardar a assinatura.
    \item No modo de verificação: fornecer ficheiro original, ficheiro de assinatura, certificado do assinante, \textit{truststore} e password para o \textit{truststore}; o programa informa se a assinatura é válida e se a cadeia de certificação foi reconhecida como confiável.
\end{itemize}

\subsection{Conclusão}

A implementação apresentada garante os requisitos fundamentais solicitados:
\begin{itemize}
    \item Assinatura e verificação de ficheiros de qualquer dimensão, graças ao processamento incremental;
    \item Suporte a funções de hash adequadas para assinatura com RSA;
    \item Validação da cadeia de certificação antes da verificação da assinatura, assegurando confiança na origem.
    \item Seleção automática da chave privada no keystore, simplificando a utilização sem necessidade de fornecer um alias.
\end{itemize}